% Lösungen Einheit 2 von 3 – Spiegelebene und Spiegelgerade bestimmen (zusammenbau v0.1)
\einheitenkopf{Einheit 2 von 3 · Spiegelebene und Spiegelgerade bestimmen}

\erg{13}{a) die Spiegelfläche gesucht – rückwärts \quad b) $\vec{n} = (1 | 2 | 0)$, $M(2 | 2 | 3)$; $E: x_1 + 2x_2 = 6$ \quad c) $S$ bleibt fest; $(4 | 4 | 2)$ auf $g$ hat das Bild $(4 | 4 | -2)$, also $h: \vec{x} = (2 | 3 | 0) + r \cdot (2 | 1 | -2)$ \quad d) beide Richtungsvektoren haben die Länge $3$; die Summe $(3 | 4 | 3)$ ist Normalenvektor: $E: 3x_1 + 4x_2 + 3x_3 = 19$ \quad e) $\overrightarrow{CW} = (-2 | 0 | -2)$, also $\vec{n} = (1 | 0 | 1)$; $U$ in $H$: $H: x_1 + x_3 = 4$ \quad f) $h$ halbiert den Winkel bei $P$. Raute aus $\overrightarrow{PB}$ (Länge $2$) und $\overrightarrow{PA}$, auf Länge $2$ gebracht: $\vec{p} + (\vec{b} - \vec{p}) + \tfrac{|\vec{b} - \vec{p}|}{|\vec{a} - \vec{p}|} \cdot (\vec{a} - \vec{p})$, hier der Punkt $(4{,}2 | 2{,}6)$ – nicht $\vec{a} + \vec{b}$}
\erg{14}{a) Fehler: $P$ statt des Mittelpunkts eingesetzt. Richtig: $M(3 | 2 | 4)$, also $E: x_1 - x_2 + x_3 = 5$ \quad b) Beim Spiegeln steht die Verbindung von Punkt und Bild senkrecht auf der Spiegelfläche – ein Vektor senkrecht zur Ebene ist ihr Normalenvektor.}
